\PassOptionsToPackage{table}{xcolor}
\documentclass[10pt,aspectratio=169]{beamer}
\newcommand{\LectureNumber}{30}
\input{course-theme.tex}
\title{Test harnesses}
\subtitle{CSC103 Programming Fundamentals / Lecture 30}
\author{Dr. Muhammad Farhan}
\institute{Associate Professor of Computer Science\\Department of Computer Science\\COMSATS University Islamabad\\Sahiwal Campus}
\date{Fall 2026}
\begin{document}
\begin{frame}\titlepage\end{frame}

\begin{frame}[fragile,t]{The problem for this lecture}
\begin{columns}[T,onlytextwidth]\begin{column}{0.60\textwidth}
Create a harness for absolute value over -5, 0 and 5. State what your method does for Integer.MIN\_VALUE.
\end{column}\begin{column}{0.35\textwidth}\small
Store manually calculated expected results in long values.
\end{column}\end{columns}
\note{Introduce the target problem before explaining the tools. Revisit it in the guided solution later.\par Syllabus reading: Reference material. CLO-4.}
\end{frame}

\begin{frame}[fragile,t]{Learning objectives}
\begin{columns}[T,onlytextwidth]\begin{column}{0.60\textwidth}
\begin{itemize}
\item Build a repeatable test harness
\item Choose expected results independently
\item Report passes and failures meaningfully
\end{itemize}
\end{column}\begin{column}{0.35\textwidth}\small
CLO-4

Tests and oracles\par A simple harness\par Equivalence and boundaries\par Repeatability and regression
\end{column}\end{columns}
\note{Explain how each objective supports the opening problem. The final questions check these objectives.\par Syllabus reading: Reference material. CLO-4.}
\end{frame}

\begin{frame}[fragile,t]{Tests and oracles}
\begin{itemize}
\item A test supplies input and checks observable behaviour.
\item An oracle specifies the expected outcome.
\item Calculate expected results independently of the implementation.
\end{itemize}
\vspace{1mm}\begin{block}{Example}
\begin{lstlisting}
Function: max(a,b)
Input: (-4,-2)
Expected: -2
Observed: result returned by the method
\end{lstlisting}
\end{block}
\note{Copying the implementation into the expected-value calculation can repeat the same defect and make a broken program appear correct.\par Syllabus reading: Reference material. CLO-4.}
\end{frame}

\begin{frame}[fragile,t]{A simple harness}
\begin{itemize}
\item A harness calls the target repeatedly and reports results.
\item Each failure should identify its case and expected value.
\item A summary helps detect missed or skipped cases.
\end{itemize}
\vspace{1mm}\begin{block}{Example}
\begin{lstlisting}
Case: negative pair
Expected: -2
Actual: 0
Status: FAIL
\end{lstlisting}
\end{block}
\note{A console harness works before a testing framework is introduced. Make the comparison and result reporting separate from the target calculation.\par Syllabus reading: Reference material. CLO-4.}
\end{frame}

\begin{frame}[fragile,t]{Equivalence and boundaries}
\begin{itemize}
\item Group inputs that should behave similarly.
\item Test transitions between groups and exceptional inputs.
\item Include empty or zero-sized input where the contract permits it.
\end{itemize}
\vspace{1mm}\begin{block}{Example}
\begin{lstlisting}
For isPass(mark):
49, 50, 51
For a valid-mark contract:
-1, 0, 100, 101
\end{lstlisting}
\end{block}
\note{These serve different purposes: threshold behaviour versus allowed-input validation. Do not assume a predicate validates its entire input domain unless its contract says so.\par Syllabus reading: Reference material. CLO-4.}
\end{frame}

\begin{frame}[fragile,t]{Repeatability and regression}
\begin{itemize}
\item Tests should produce the same result under the same conditions.
\item Control random seeds, data files and other changing dependencies.
\item Keep a small test that exposes each fixed bug.
\end{itemize}
\vspace{1mm}\begin{block}{Example}
\begin{lstlisting}
Regression case:
Input: two negative numbers
Expected: larger negative number
Purpose: detect zero initialisation bug
\end{lstlisting}
\end{block}
\note{A passing test suite supplies evidence about selected cases. It does not prove that every input is correct.\par Syllabus reading: Reference material. CLO-4.}
\end{frame}

\begin{frame}[fragile,t]{The expected result must be independent}
\begin{itemize}
\item A harness compares actual behaviour to an oracle derived from the specification.
\item Using the same potentially faulty implementation on both sides proves little.
\item Literal expected results work well for a small set of manually solved cases.
\end{itemize}
\vspace{1mm}\begin{block}{Example}
\begin{lstlisting}
Input -5 has expected absolute value 5.
The expected value comes from the mathematical definition.
\end{lstlisting}
\end{block}
\note{Explain the mechanism using the displayed example. A harness compares actual behaviour to an oracle derived from the specification. Using the same potentially faulty implementation on both sides proves little. Literal expected results work well for a small set of manually solved cases.\par Syllabus reading: Reference material. CLO-4.}
\end{frame}

\begin{frame}[fragile,t]{Numeric limits belong in the contract}
\begin{itemize}
\item Integer.MIN\_VALUE has no positive counterpart representable as int.
\item Returning long lets an absolute-value function represent that magnitude.
\item The conversion must happen before negation to avoid int overflow.
\end{itemize}
\vspace{1mm}\begin{block}{Example}
\begin{lstlisting}
long widened = value;
return widened < 0 ? -widened : widened;
\end{lstlisting}
\end{block}
\note{Explain the mechanism using the displayed example. Integer.MIN\_VALUE has no positive counterpart representable as int. Returning long lets an absolute-value function represent that magnitude. The conversion must happen before negation to avoid int overflow.\par Syllabus reading: Reference material. CLO-4.}
\end{frame}


\begin{frame}[fragile,t]{A test harness compares two independent paths}
\begin{center}\begin{tikzpicture}
\node[box] (a) at (0,0) {Chosen input};
\node[alt] (b) at (4,1) {Program result};
\node[hot] (c) at (4,-1) {Expected result};
\node[box] (d) at (8,0) {Compare and report};
\draw[arr](a.east)--++(.6,0)|-(b.west);\draw[arr](a.east)--++(.6,0)|-(c.west);\draw[arr](b.east)--++(.6,0)|-(d.west);\draw[arr](c.east)--++(.6,0)|-(d.west);
\end{tikzpicture}\end{center}
\begin{block}{What to notice}The expected result must come from the specification, not from calling the same implementation again.\end{block}
\note{Trace each labelled connection. The expected result must come from the specification, not from calling the same implementation again.}
\end{frame}
\begin{frame}[fragile,t]{Key distinctions}
\small\renewcommand{\arraystretch}{1.5}
\rowcolors{2}{pale}{white}
\begin{tabularx}{\textwidth}{>{\raggedright\arraybackslash}X>{\raggedright\arraybackslash}X>{\raggedright\arraybackslash}X}
\rowcolor{navy}
\color{white}\bfseries Input & \color{white}\bfseries Expected long result & \color{white}\bfseries Reason for case\\
-5 & 5 & Negative ordinary input\\
0 & 0 & Zero boundary\\
5 & 5 & Positive ordinary input\\
Integer.MIN\_VALUE & 2147483648 & Minimum int magnitude\\
\end{tabularx}
\vspace{3mm}\footnotesize These distinctions determine which operation or control structure fits the problem.
\note{\par Syllabus reading: Reference material. CLO-4.}
\end{frame}

\begin{frame}[fragile,t]{First example: methods}
\begin{lstlisting}
static int max(int a, int b) { return a >= b ? a : b; }
\end{lstlisting}
\note{Method definitions belong inside the class and outside main. Explain the parameters and return values.\par Syllabus reading: Reference material. CLO-4.}
\end{frame}

\begin{frame}[fragile,t]{First example: execution}
\begin{lstlisting}
int[][] cases = {{2, 5, 5}, {-4, -2, -2}, {3, 3, 3}};
int passed = 0;
for (int[] c : cases) {
  int actual = max(c[0], c[1]);
  if (actual == c[2]) passed++;
  else System.out.println("FAIL expected " + c[2] + " got " + actual);
}
System.out.println("Passed " + passed + "/" + cases.length);
\end{lstlisting}
\note{This is the main body from Java\_Examples/Lecture30.java. The source file includes the complete class. Predict the result before the next slide.\par Syllabus reading: Reference material. CLO-4.}
\end{frame}

\begin{frame}[fragile,t]{First example: result}
\begin{columns}[T,onlytextwidth]\begin{column}{0.60\textwidth}
Passed 3/3\par Each row contains first input, second input and expected result.\par The equal-input case checks ties.
\end{column}\begin{column}{0.35\textwidth}\small
Equivalence and boundaries

Repeatability and regression
\end{column}\end{columns}
\note{Expected output and interpretation: Passed 3/3
Each row contains first input, second input and expected result.
The equal-input case checks ties.\par Syllabus reading: Reference material. CLO-4.}
\end{frame}

\begin{frame}[fragile,t]{Guided practice}
\begin{columns}[T,onlytextwidth]\begin{column}{0.60\textwidth}
Create a harness for absolute value over -5, 0 and 5. State what your method does for Integer.MIN\_VALUE.
\end{column}\begin{column}{0.35\textwidth}\small
Attempt the solution before continuing.

Use the definitions and examples from this lecture.
\end{column}\end{columns}
\note{Give students 8 minutes to attempt the problem. The following slides provide a complete solution. Original solution guidance: Expected ordinary results: 5,0,5. Since the positive counterpart of Integer.MIN\_VALUE does not fit int, explicitly reject it, use long, or document the behaviour. Do not invent an in-range result.\par Syllabus reading: Reference material. CLO-4.}
\end{frame}

\begin{frame}[fragile,t]{Solution strategy}
\begin{columns}[T,onlytextwidth]\begin{column}{0.60\textwidth}
\begin{itemize}
\item Store manually calculated expected results in long values.
\item Call the absolute-value method for each input.
\item Print expected and actual values with a clear pass or fail label.
\end{itemize}
\end{column}\begin{column}{0.35\textwidth}\small
The next program uses fixed inputs so its result can be reproduced.
\end{column}\end{columns}
\note{Discuss why each step is needed. State the input assumptions before executing the program.\par Syllabus reading: Reference material. CLO-4.}
\end{frame}

\begin{frame}[fragile,t]{Complete solution (1/2)}
\begin{lstlisting}
public class Solution30 {
  public static void main(String[] args) throws Exception {
    int[] inputs = {-5, 0, 5, Integer.MIN_VALUE};
    long[] expected = {5L, 0L, 5L, 2147483648L};
    for (int i = 0; i < inputs.length; i++) {
      long actual = absolute(inputs[i]);
      String status = actual == expected[i] ? "PASS" : "FAIL";
      System.out.println(status + " " + inputs[i] + " => " + actual);
\end{lstlisting}
\note{Complete runnable file: Java\_Examples/Solution30.java. Inputs are fixed in the program.  \par Syllabus reading: Reference material. CLO-4.}
\end{frame}

\begin{frame}[fragile,t]{Complete solution (2/2)}
\begin{lstlisting}
    }
  }
  static long absolute(int value) {
    long widened = value;
    return widened < 0 ? -widened : widened;
  }
}
\end{lstlisting}
\note{Complete runnable file: Java\_Examples/Solution30.java. Inputs are fixed in the program.  \par Syllabus reading: Reference material. CLO-4.}
\end{frame}

\begin{frame}[fragile,t]{Execution trace}
\small\renewcommand{\arraystretch}{1.5}
\rowcolors{2}{pale}{white}
\begin{tabularx}{\textwidth}{>{\raggedright\arraybackslash}X>{\raggedright\arraybackslash}X>{\raggedright\arraybackslash}X>{\raggedright\arraybackslash}X}
\rowcolor{navy}
\color{white}\bfseries Input & \color{white}\bfseries Widened value & \color{white}\bfseries Negation needed? & \color{white}\bfseries Actual result\\
-5 & -5L & Yes & 5L\\
0 & 0L & No & 0L\\
5 & 5L & No & 5L\\
Minimum int & -2147483648L & Yes & 2147483648L\\
\end{tabularx}
\vspace{3mm}\footnotesize Follow the changing state in execution order.
\note{Manually derived trace for the complete solution. Store manually calculated expected results in long values. Call the absolute-value method for each input. Print expected and actual values with a clear pass or fail label.\par Syllabus reading: Reference material. CLO-4.}
\end{frame}

\begin{frame}[fragile,t]{Solution result}
\begin{columns}[T,onlytextwidth]\begin{column}{0.60\textwidth}
PASS -5 =\textgreater{} 5\par PASS 0 =\textgreater{} 0\par PASS 5 =\textgreater{} 5\par PASS -2147483648 =\textgreater{} 2147483648
\end{column}\begin{column}{0.35\textwidth}\small
Why this result follows

Print expected and actual values with a clear pass or fail label.
\end{column}\end{columns}
\note{Expected result: PASS -5 =\textgreater{} 5
PASS 0 =\textgreater{} 0
PASS 5 =\textgreater{} 5
PASS -2147483648 =\textgreater{} 2147483648\par Syllabus reading: Reference material. CLO-4.}
\end{frame}

\begin{frame}[fragile,t]{A common incorrect approach}
int expected = buggyMax(a,b);\par int actual = buggyMax(a,b);
\vspace{1mm}\begin{block}{Example}
\begin{lstlisting}
The oracle repeats the same implementation. Write the expected value from the contract and a manually solved case.
\end{lstlisting}
\end{block}
\note{Ask students to predict the failure before discussing the explanation. The right side gives the correction.\par Syllabus reading: Reference material. CLO-4.}
\end{frame}

\begin{frame}[fragile,t]{Boundary and failure checks}
\begin{columns}[T,onlytextwidth]\begin{column}{0.60\textwidth}
Check the stated input assumptions.

Build a harness with at least six cases for your average or maximum method and retain one deliberate bug demonstration.
\end{column}\begin{column}{0.35\textwidth}\small
Expected ordinary results: 5,0,5. Since the positive counterpart of Integer.MIN\_VALUE does not fit int, explicitly reject it, use long, or document the behaviour. Do not invent an in-range result.
\end{column}\end{columns}
\note{Use the specification to derive each expected result. Exercise solution: Expected ordinary results: 5,0,5. Since the positive counterpart of Integer.MIN\_VALUE does not fit int, explicitly reject it, use long, or document the behaviour. Do not invent an in-range result.\par Syllabus reading: Reference material. CLO-4.}
\end{frame}

\begin{frame}[fragile,t]{Check your understanding}
\begin{columns}[T,onlytextwidth]\begin{column}{0.60\textwidth}
What makes a failed test useful? Why include a tie case for maximum?
\end{column}\begin{column}{0.35\textwidth}\small
Write a prediction and explain the rule that supports it.
\end{column}\end{columns}
\note{Answer: It names the case and reports expected versus actual. Ties exercise equality rather than only strict ordering.\par Syllabus reading: Reference material. CLO-4.}
\end{frame}

\begin{frame}[fragile,t]{Answer and explanation}
\begin{columns}[T,onlytextwidth]\begin{column}{0.60\textwidth}
It names the case and reports expected versus actual. Ties exercise equality rather than only strict ordering.
\end{column}\begin{column}{0.35\textwidth}\small
The oracle repeats the same implementation. Write the expected value from the contract and a manually solved case.
\end{column}\end{columns}
\note{Reconnect the answer to the relevant definition rather than treating it as a fact to memorise.\par Syllabus reading: Reference material. CLO-4.}
\end{frame}


\begin{frame}[fragile,t]{Compare cases and predict outcomes}
\small\renewcommand{\arraystretch}{1.5}\rowcolors{2}{pale}{white}
\begin{tabularx}{\textwidth}{>{\raggedright\arraybackslash}X>{\raggedright\arraybackslash}X>{\raggedright\arraybackslash}X}\rowcolor{navy}
\color{white}\bfseries Test & \color{white}\bfseries Expected max & \color{white}\bfseries Fault it can expose\\
max(3,7) & 7 & Wrong comparison\\
max(-3,-7) & -3 & Incorrect zero initialisation\\
max(5,5) & 5 & Mishandled equality\\
\end{tabularx}\par\vspace{3mm}\begin{exampleblock}{Discuss}Explain each expected result before running the example. Which case would reveal a wrong assumption?\end{exampleblock}
\note{Read the first two columns and invite predictions before discussing the outcome.}
\end{frame}

\begin{frame}[fragile,t]{Apply the idea}
\begin{alertblock}{Independent practice}Build a small explicit check for max(-3,-7). Explain why Java assert statements can be an unreliable default harness without the required runtime option.\end{alertblock}\vspace{3mm}\begin{enumerate}\item Work through the input by hand.\item Write or correct the program.\item Justify the expected result.\end{enumerate}\note{Allow 3–5 minutes, then compare answers in pairs. Build a small explicit check for max(-3,-7). Explain why Java assert statements can be an unreliable default harness without the required runtime option.}
\end{frame}

\begin{frame}[fragile,t]{Solution and reasoning}
\begin{lstlisting}
int actual = Math.max(-3, -7);
int expected = -3;
if (actual != expected) {
  throw new AssertionError("expected -3, got " + actual);
}
System.out.println("Passed");
\end{lstlisting}\vspace{1mm}\small\begin{itemize}
\item The expected value is reasoned from the ordering of negative numbers.
\item The explicit check executes normally.
\item Java language assertions are disabled by default unless enabled, for example with -ea.
\end{itemize}\note{Answer: The expected value is reasoned from the ordering of negative numbers. The explicit check executes normally. Java language assertions are disabled by default unless enabled, for example with -ea.}
\end{frame}

\begin{frame}[fragile,t]{A harness for the digit{-}sum exercise}
\begin{itemize}
\item Turn the source example into a method contract with ordinary and boundary cases.
\item An expected value must be derived independently of the method under test.
\item A harness should fail visibly when a comparison fails.
\end{itemize}
\vspace{3mm}\footnotesize\textcolor{accent}{Source: supplied ISB campus slides. See speaker notes and ISB\_Source\_Map.md.}
\note{Adapted from the supplied ISB campus folder: Methods Practice.pptx, slides 2, 3. Instructor{-}authored testing adaptation of the supplied digit{-}sum exercise; the source does not contain this harness. Source exercise deck credits Instructor Khurram Iqbal.}
\end{frame}

\begin{frame}[fragile,t]{Worked problem: A harness for the digit{-}sum exercise}
\begin{block}{Task}Check digit sums for 234, 0 and 1001 with explicit comparisons; report how many pass.\end{block}
\small Input: Values are fixed in the program for a reproducible demonstration.\par\vspace{3mm}\begin{exampleblock}{Before running}Predict the output and identify the variables that change.\end{exampleblock}
\note{Adapted from the supplied ISB campus folder: Methods Practice.pptx, slides 2, 3. Instructor{-}authored testing adaptation of the supplied digit{-}sum exercise; the source does not contain this harness. Source exercise deck credits Instructor Khurram Iqbal.}
\end{frame}

\begin{frame}[fragile,t]{Complete ISB example (1/2)}
\begin{lstlisting}
public class IsbLecture30 {
  public static void main(String[] args) throws Exception {
    long[] inputs = {234, 0, 1001};
    int[] expected = {9, 0, 2};
    for (int i = 0; i < inputs.length; i++) {
      int actual = sumDigits(inputs[i]);
      if (actual != expected[i]) throw new AssertionError("case " + i);
    }
    System.out.println("Passed 3/3");
\end{lstlisting}
\note{Adapted from the supplied ISB campus folder: Methods Practice.pptx, slides 2, 3. Instructor{-}authored testing adaptation of the supplied digit{-}sum exercise; the source does not contain this harness. Source exercise deck credits Instructor Khurram Iqbal. Complete file: ISB\_Java\_Examples/IsbLecture30.java.}
\end{frame}

\begin{frame}[fragile,t]{Complete ISB example (2/2)}
\begin{lstlisting}
  }
  static int sumDigits(long n) {
    if (n < 0) throw new IllegalArgumentException("nonnegative required");
    int sum = 0;
    while (n != 0) { sum += n % 10; n /= 10; }
    return sum;
  }
}
\end{lstlisting}
\note{Adapted from the supplied ISB campus folder: Methods Practice.pptx, slides 2, 3. Instructor{-}authored testing adaptation of the supplied digit{-}sum exercise; the source does not contain this harness. Source exercise deck credits Instructor Khurram Iqbal. Complete file: ISB\_Java\_Examples/IsbLecture30.java.}
\end{frame}

\begin{frame}[fragile,t]{Worked trace and expected result}
\small\renewcommand{\arraystretch}{1.4}\rowcolors{2}{pale}{white}
\begin{tabularx}{\textwidth}{>{\raggedright\arraybackslash}X>{\raggedright\arraybackslash}X>{\raggedright\arraybackslash}X}\rowcolor{navy}
\color{white}\bfseries Input & \color{white}\bfseries Expected sum & \color{white}\bfseries Reason\\
234 & 9 & Ordinary source case\\
0 & 0 & Zero boundary\\
1001 & 2 & Internal zero digits\\
\end{tabularx}\par\vspace{2mm}
\begin{block}{Expected console output}\ttfamily\footnotesize Passed 3/3\end{block}
\note{Adapted from the supplied ISB campus folder: Methods Practice.pptx, slides 2, 3. Instructor{-}authored testing adaptation of the supplied digit{-}sum exercise; the source does not contain this harness. Source exercise deck credits Instructor Khurram Iqbal. Trace and expected values independently worked for this edition.}
\end{frame}

\begin{frame}[fragile,t]{Extend the ISB example}
\begin{alertblock}{Practice}Would a method that always returns 9 pass this harness?\end{alertblock}\vspace{3mm}\begin{exampleblock}{Explain your reasoning}It would pass only the first case and fail the zero case. The tests distinguish more than the original example.\end{exampleblock}
\note{Adapted from the supplied ISB campus folder: Methods Practice.pptx, slides 2, 3. Instructor{-}authored testing adaptation of the supplied digit{-}sum exercise; the source does not contain this harness. Source exercise deck credits Instructor Khurram Iqbal. Answer: It would pass only the first case and fail the zero case. The tests distinguish more than the original example.}
\end{frame}
\begin{frame}[fragile,t]{Lecture summary}
\begin{columns}[T,onlytextwidth]\begin{column}{0.60\textwidth}
\begin{itemize}
\item a repeatable test harness
\item expected results independently
\item Report passes and failures meaningfully
\end{itemize}
\end{column}\begin{column}{0.35\textwidth}\small
\begin{itemize}
\item Store manually calculated expected results in long values.
\item Call the absolute-value method for each input.
\item Print expected and actual values with a clear pass or fail label.
\end{itemize}
\end{column}\end{columns}
\note{Ask students to give one concrete example for the concepts listed. Use the worked solution to connect the topics.\par Syllabus reading: Reference material. CLO-4.}
\end{frame}

\begin{frame}[fragile,t]{After-class practice}
\begin{columns}[T,onlytextwidth]\begin{column}{0.60\textwidth}
Build a harness with at least six cases for your average or maximum method and retain one deliberate bug demonstration.
\end{column}\begin{column}{0.35\textwidth}\small
Syllabus reading\par Reference material

Next lecture: Unit testing with JUnit
\end{column}\end{columns}
\note{Reading labels follow the supplied outline. Check topic headings against the available textbook edition. Require a program and expected results for the practice task.\par Syllabus reading: Reference material. CLO-4.}
\end{frame}
\end{document}
